%========================================================================================
%   Chapter 8 — Challenges & Solutions
%========================================================================================
\chapter{Challenges \& Solutions}

This chapter highlights the main technical and operational challenges encountered during the project lifecycle and the strategies adopted to overcome them.  
Each problem is presented with its root cause, corrective actions, and impact on the final system’s quality and security posture.

%------------------------------------------------
\section{Development Challenges}

\begin{itemize}
    \item \textbf{JWT Authentication Conflicts:}  
    Early Spring Security configurations introduced cyclic dependencies between authentication beans and user details loading.  
    \textbf{Solution:} Implemented a custom \texttt{SimpleUserDetailsService} leveraging BCrypt password hashing and refactored beans to load through environment-based credentials, breaking circular injection chains.

    \item \textbf{Foreign Key Constraint Errors:}  
    CRUD operations involving \texttt{Client}, \texttt{Driver}, and \texttt{TransportLog} entities failed due to cascading deletions.  
    \textbf{Solution:} Annotated transactional service methods with \texttt{@Transactional} and added explicit entity loading to preserve relational integrity.

    \item \textbf{React Authentication State Handling:}  
    Token expiration caused inconsistent authentication states between components.  
    \textbf{Solution:} Introduced a centralized \texttt{AuthContext} and Axios interceptors to refresh tokens automatically and perform clean logout on 401 responses.

    \item \textbf{DTO–Entity Mapping Complexity:}  
    Manual mapping generated redundant code.  
    \textbf{Solution:} Adopted \texttt{MapStruct} to automate DTO conversions, improving maintainability.
\end{itemize}

%------------------------------------------------
\section{Infrastructure \& Deployment Challenges}

\begin{itemize}
    \item \textbf{Container Startup Dependencies:}  
    Spring Boot containers started before PostgreSQL initialization, causing connection errors.  
    \textbf{Solution:} Added Docker health checks and conditional startup logic:
\begin{verbatim}
db:
  healthcheck:
    test: ["CMD-SHELL", "pg_isready -U ${POSTGRES_USER}"]
app:
  depends_on:
    db:
      condition: service_healthy
\end{verbatim}

    \item \textbf{Large Image Sizes:}  
    Dependencies inflated container size up to 500 MB.  
    \textbf{Solution:} Multi-stage Docker builds reduced final image size by 60\,\%, keeping runtime layers minimal and secure.

    \item \textbf{Environment Configuration Management:}  
    Managing separate credentials for staging and production led to inconsistencies.  
    \textbf{Solution:} Introduced environment-specific Spring profiles and Docker-injected environment variables for JWT, database URLs, and secrets.

    \item \textbf{Terraform State Conflicts:}  
    Parallel executions created drift between environments.  
    \textbf{Solution:} Migrated Terraform state to a remote backend with locking enabled to guarantee atomic operations.
\end{itemize}

%------------------------------------------------
\section{Security Challenges}

\begin{itemize}
    \item \textbf{Secret Management:}  
    Early prototypes used plaintext JWT secrets and credentials.  
    \textbf{Solution:} Moved all secrets to environment variables and encrypted vault files managed by Ansible Vault (\texttt{group\_vars/staging/vault.yml}).

    \item \textbf{API Security Testing:}  
    Risk of leaking credentials during automated tests.  
    \textbf{Solution:} Created a dedicated \texttt{TestSecurityConfig} with mock authentication providers, isolating tests from real secrets.

    \item \textbf{Validation Error Disclosure:}  
    Raw exceptions exposed internal details to clients.  
    \textbf{Solution:} Added a global \texttt{@ControllerAdvice} handler to sanitize responses and standardize error formats.

    \item \textbf{Firewall Misconfiguration:}  
    Backend ports were unintentionally accessible publicly during initial tests.  
    \textbf{Solution:} Re-applied UFW and DigitalOcean Cloud Firewall rules, restricting API access to frontend and bastion hosts only.
\end{itemize}

%------------------------------------------------
\section{Testing \& CI/CD Challenges}

\begin{itemize}
    \item \textbf{Database Testing Isolation:}  
    Integration tests risked mutating production data.  
    \textbf{Solution:} Configured tests to use an in-memory H2 database and mocked repositories with Mockito for pure unit testing.

    \item \textbf{Security Flow Testing:}  
    Complex JWT-based filters required full authentication simulation.  
    \textbf{Solution:} Implemented \texttt{MockMvc} integration tests to validate access rules and token lifecycles.

    \item \textbf{CI/CD Workflow Failures:}  
    Initial GitHub Actions failed on Terraform or Ansible steps due to secret misconfiguration.  
    \textbf{Solution:} Segmented workflows (build/test, scan, deploy) with environment-specific secrets and conditional job triggers for staging/production.
\end{itemize}

%------------------------------------------------
\section{Performance Optimization Solutions}

\begin{itemize}
    \item \textbf{API Optimization:}  
    Implemented pagination and projection across all endpoints, reducing payload sizes.  

    \item \textbf{Query Tuning:}  
    Optimized JPA queries and indexed primary keys, cutting average DB response to 67 ms.  

    \item \textbf{Startup Speed:}  
    Cached Maven dependencies and precompiled layers, reducing container startup to 18 s.  

    \item \textbf{Static Resource Caching:}  
    Configured Nginx caching headers, improving page-load time by 45\,\%.
\end{itemize}

%------------------------------------------------
\section{Organizational \& Workflow Challenges}

\begin{itemize}
    \item \textbf{Collaboration and Version Control:}  
    Merging Terraform and Ansible updates between developers introduced conflicts.  
    \textbf{Solution:} Adopted Git branching strategy (\texttt{features/back}, \texttt{features/front}, \texttt{infra/}) with pull request reviews.

    \item \textbf{Documentation Synchronization:}  
    Rapid infrastructure iterations caused outdated diagrams.  
    \textbf{Solution:} Automated regeneration of PlantUML diagrams and inventory files through Makefile targets (\texttt{make docs}).

    \item \textbf{Remote Deployment Connectivity:}  
    SSH agent forwarding failed in early Ansible executions.  
    \textbf{Solution:} Enabled ProxyJump through bastion and verified agent forwarding using \texttt{ssh -T git@github.com} tests before deployment.
\end{itemize}

%------------------------------------------------
\section{Key Results and Impact}

\begin{table}[H]
\centering
\begin{tabular}{|l|c|c|}
\hline
\textbf{Metric} & \textbf{Before} & \textbf{After} \\ \hline
Container Image Size & 500 MB & 200 MB \\
API Response Time & 800 ms & 245 ms \\
Security Vulnerabilities & 12 issues & 0 critical \\
Test Coverage & 60 \% & 95 \% \\
Deployment Time & 15 min & 8 min \\ \hline
\end{tabular}
\caption{Resolution Outcomes and Performance Improvements}
\end{table}

%------------------------------------------------
\section{Summary of Lessons Learned}

The project’s challenges revealed the value of \textbf{DevSecOps integration} throughout the lifecycle:
\begin{itemize}
    \item Early IaC adoption prevented configuration drift across environments.  
    \item Security-by-design minimized vulnerabilities without slowing delivery.  
    \item Automated CI/CD pipelines increased reliability and developer confidence.  
    \item Documentation and modularity accelerated debugging and maintenance.  
\end{itemize}

The systematic resolution of these challenges led to a robust, secure, and reproducible system architecture that can scale easily and serve as a blueprint for future DevSecOps projects.
